\section{Policy gradients and the likelihood-ratio contract}
\label{sec:C08:policygradient}
For a trajectory law $p_\theta(\tau)$ and a parameter-independent return $R(\tau)$, the score-function identity gives $\nabla_\theta J=\mathbb E[R(\tau)\nabla_\theta\log p_\theta(\tau)]$ under differentiation and integrability assumptions. Fixed environment-transition factors cancel from the score. All parameterized agent decisions remain, including memory writes, context edits, routing, or stopping when these are learned actions. Tool observations themselves are not policy-generated action tokens. A state-conditional baseline independent of the sampled action can reduce variance without changing the expected gradient; an action-dependent baseline requires correction.

For example, let a writer sample memory updates from $\omega_\psi(m_{t+1}\mid H_t,M_t)$ and let a separate actor use $\pi_\theta(a_t\mid C_t)$, with context a deterministic function of the realized history and memory. Assume $\psi$ enters the trajectory law only through the writer. Then
\begin{equation}
 \nabla_\psi J=
 \mathbb E\!\left[R(\tau)
 \sum_t\nabla_\psi\log
 \omega_\psi(m_{t+1}\mid H_t,M_t)\right].
\end{equation}
Training only the actor's action-token loss omits this writer gradient. It can still define a deliberate alternating or frozen-writer procedure, but is not joint optimization of both components. Shared parameters require adding every path through which they affect the sampling law; a differentiable parameter-dependent reward adds its direct derivative. This decomposition gives a concrete audit target for harness-native training and the co-adaptation analysis in Chapter~\ref{ch:C17}.

PPO uses a surrogate based on the ratio $r_t(\theta)=\pi_\theta(a_t\mid C_t)/\pi_{\mathrm{old}}(a_t\mid C_t)$ and an advantage estimate $\widehat A_t$. Its clipped objective is
\begin{equation}
 \mathbb E\min\{r_t\widehat A_t,
 \operatorname{clip}(r_t,1-\epsilon,1+\epsilon)\widehat A_t\}.
\end{equation}
The practical clipped method is not a universal monotonic-improvement guarantee. It relies on sampled data, estimated advantages, optimization choices, and an appropriate proximity regime~\cite{schulman2017ppo}. In long-horizon agents, sparse rewards, variable action lengths, and tool delays make advantage estimation and sample accounting especially consequential.

The denominator must describe the behavior that generated the action under the same conditioning information. If a trace is later compacted and the numerator is evaluated on the final summary while the denominator comes from the live pre-compaction context, the ratio is not the intended conditional likelihood ratio. Likewise, a grammar mask or temperature used during sampling changes behavior probabilities. Storing base-model logits without the sampling transformation can produce an incorrect correction. The trace should retain the live context identity, sampled token probabilities after relevant transforms, and the exact action mask.

A trajectory-level importance ratio is the product of conditional ratios over its sampled actions. Its variance can grow severely with horizon. Token-level or per-decision approximations can reduce variance while changing the estimator and its assumptions. Clipping further introduces bias deliberately. These are not reasons to reject the methods; they are reasons to state which expectation is being estimated. An agent result should distinguish an exact identity, a biased practical estimator, and an empirical observation that the estimator worked in the reported setting.

